% Artifact appendix aligned with the repository README.
\section{Artifact Evaluation}
\label{sec:artifact}

\subsection{Overview and Access}

The artifact provides SME and Harmony implementations for GPT-2
training, including model configurations, generated layer code,
profiles, four-GPU schedules, and Bash/Python experiment launchers.
It supports comparison of memory consumption, throughput, and
maximum trainable model size using pipeline parallelism (PP) and
data parallelism (DP). Source code and the complete reproduction
instructions are available at
\url{https://github.com/dinwenu/temprepo}.
The top-level README is the entry point for installation, dataset
preparation, script selection, and expected results.

\subsection{Artifact Checklist}

\begin{itemize}
  \item \textbf{Algorithm and programs:} Task-level parameter and
  gradient memory management with layer-granularity NVMe parameter
  swapping; SME and Harmony GPT-2 workloads.
  \item \textbf{Models:} 1.5B, 2.6B, 5.9B, 8.0B, and 10.5B GPT-2
  variants; sequence length 1024, FP32 precision, and activation
  recomputation. Configurations and model code are included.
  Pretrained weights are not bundled; default runs initialize
  weights from the configurations.
  \item \textbf{Data:} WikiText-103, downloaded separately from the
  \texttt{wikitext-103-v1} subset of Salesforce's Hugging Face
  repository. Downloads total approximately 313\,MB; restored text
  totals approximately 542\,MB. The GPT-2 tokenizer is downloaded
  separately on first use.
  \item \textbf{Software:} Linux x86-64, Python
  3.9.19, PyTorch 1.12.1, CUDA 11.3, and DeepSpeed 0.11.1; the paper
  reports NCCL 2.10. Dependencies are specified in
  \texttt{environment.yml}. A compatible NVIDIA driver, C++ and
  CUDA compilers, Ninja, and libaio development files are required.
  \item \textbf{Hardware and storage:} The reference server has
  four RTX 3090 GPUs (24\,GB each), two 12-core Intel Xeon 4214R
  CPUs, 256\,GB DRAM, a 2\,TB Samsung PM981 NVMe SSD, and PCIe
  3.0$\times$16 GPU connections. At least \textbf{150\,GB of free
  SSD space} is recommended.
  \item \textbf{Execution and run-time state:} Run experiments
  sequentially on four GPUs. Throughput averages exclude the first
  half of completed iterations as warm-up. The 256\,GB capacity
  experiment runs directly on the full-memory server; the 128\,GB
  experiment uses NUMA binding or a collectively limited no-swap
  cgroup.
  \item \textbf{Metrics and output:} CPU/GPU memory usage, iteration
  time, throughput (samples/s), training loss, and maximum trainable
  model size. Console logs, loss arrays, and \path{command.json}
  are produced. Capacity sweeps also save case logs,
  \path{constraints.json}, and \path{summary.csv}.
\end{itemize}

\subsection{Installation, Data, and Basic Test}

From the repository root, create the environment and configure a
writable directory on the NVMe SSD:
\begin{verbatim}
conda env create -n sme -f environment.yml
conda activate sme
export SME_NVME_PATH=/path/on/nvme/sme
mkdir -p "$SME_NVME_PATH"
ulimit -n 51200
\end{verbatim}
The shell's hard limit must permit this open-file limit. Install
libaio development files and set library/include paths as described
in the README. DeepSpeed builds its I/O extension on first use;
the custom swapper is imported from the repository, without copying
files into the installed DeepSpeed package.

Follow the README's download and conversion commands for
\texttt{wikitext-103-v1} at the pinned Hugging Face revision
specified there.
Preserve row and shard order; restore each empty text row to one
space followed by a newline, then concatenate and write UTF-8 bytes.
Do not strip whitespace or add newlines to nonempty rows. The
provided converter checks all three output files against the
original experiment files' SHA256 values. The training file must be
\path{data/wikitext-103-tokens/wiki.train.tokens}.

All remaining commands below run from the launcher directory:
\begin{verbatim}
cd harmony/4_runtime/gpt2_huggingface
bash run_gpt2_work5_2-6b.sh --check
bash run_gpt2_work5_xl.sh --num-iters 2
\end{verbatim}
The input check prints that packaged inputs are present; it does
not exercise training. The two-iteration test should initialize
four ranks, print finite loss and positive throughput, finish with
an average-throughput and memory summary, and exit successfully.

\subsection{Experiments and Expected Results}

\paragraph{Memory and throughput.}
Standard schedules cover global minibatch sizes 16, 32, 64, 128,
and 256. For example, run SME-DP with the 5.9B model:
\begin{verbatim}
bash run_gpt2_work5_5-9b.sh \
  --schedule D32_vDP_N4_Top1 --num-iters 10
\end{verbatim}
Use \texttt{vPP} for PP. The corresponding Harmony launcher is
\texttt{run\_gpt2\_5-9b.sh}; the README lists launchers for all five
sizes. Each method should use matching model, data, schedule, and
precision settings. Historical internal names such as
\texttt{gpt2\_5b} identify serialized artifacts; public script names
use the corrected sizes.

Each iteration reports loss and throughput. Rank~0 prints
\texttt{Avg Iter Time} and \texttt{Avg Throughput}; average
throughput equals global minibatch size divided by the mean
iteration duration after warm-up. Save console output with
\texttt{tee}. The reported CPU \texttt{occupied} values are
system-wide \texttt{total - available} snapshots; peak-memory
comparisons require consistent sampling over the full execution.

\paragraph{Maximum trainable model size (Section 4.2).}
Verify the dedicated schedules and run the capacity sweeps:
\begin{verbatim}
python3 prepare_capacity_schedules.py --check
bash run_section4_2_128gb.sh
bash run_section4_2_256gb.sh
\end{verbatim}
These schedules use global minibatch size 4, four GPUs, forward
and backward microbatch size 1, and one vLayer per pack. Each sweep
covers SME/Harmony, five model sizes, and PP/DP, for 20 cases with
four iterations per case. The 256\,GB sweep runs directly on the
reference server. For 128\,GB, the default
\texttt{--cpunodebind=0 --membind=0} is appropriate only when node~0
has approximately 128\,GB; binding alone does not set a byte limit.
The README describes node overrides and the cgroup alternative.
Adding \texttt{--check} to a sweep inspects inputs and reports memory
constraints without starting training.

On the reference platform, the expected largest supported models
for both PP and DP are:
\begin{center}
\begin{tabular}{lrr}
\hline
CPU memory & SME & Harmony \\
\hline
128\,GB & 5.9B & 1.5B \\
256\,GB & 10.5B & 2.6B \\
\hline
\end{tabular}
\end{center}
SME should use less CPU memory for comparable workloads. Inspect
case logs and completion/failure status in \texttt{summary.csv};
record OOM separately from other failures. Absolute timing and
throughput vary with hardware and system load. No numerical
variation tolerance has been empirically established.
